% !TeX program = lualatex
\documentclass[11pt,a4paper]{article}

% Definit le chemin vers les ressources LaTeX du projet
\makeatletter
\def\input@path{{../../../tex/}}
\makeatother
\usepackage{college-eval-v2}

% Métadonnées du contrôle
\examsetup{6\textsuperscript{e}}{Mathématiques}{Géométrie et Division euclidienne}{2}{50 minutes}{Calculatrice autorisée}

\begin{document}

% Filigrane et bannière
\ExamBanner
\StudentInfoBox

\smallspace

% === Barème avec nouvelle fonction ===
\bareme{
\exbadge{1}{5} \quad
\exbadge{2}{4} \quad
\exbadge{3}{3} \quad
\exbadge{4}{4} \quad
\exbadge{5}{4}
}

% === Consignes ===
\begin{consignebox}
\faIcon{info-circle} \textbf{CONSIGNES IMPORTANTES}
\begin{itemize}[leftmargin=*,noitemsep,topsep=5pt]
\item[\faIcon{check}] Répondez directement sur le sujet
\item[\faIcon{check}] Les constructions géométriques doivent être soignées et précises
\item[\faIcon{check}] Justifiez vos réponses quand c'est demandé
\item[\faIcon{pencil-alt}] Écrivez lisiblement
\end{itemize}
\end{consignebox}

\exspace

% ============================================
% EXERCICE 1 : Vocabulaire et notations (5 points)
% ============================================
\ExTemplate{1}{Vocabulaire et notations}{5}{book}{
\textbf{1.} Donne la signification (lecture) de chaque notation : \points{2,5}

\smallspace

\begin{itemize}[label=$\bullet$,leftmargin=*]
\item $[EF]$ se lit : \dotline{8cm}

\item $(MN)$ se lit : \dotline{8cm}

\item $[RS)$ se lit : \dotline{8cm}
\end{itemize}

\vspace{0.4cm}

\textbf{2.} Écris avec les notations mathématiques : \points{2,5}

\smallspace

\begin{itemize}[label=$\bullet$,leftmargin=*]
\item La droite passant par $A$ et $B$ : \dotline{4cm}

\item Le segment d'extrémités $C$ et $D$ : \dotline{4cm}

\item La demi-droite d'origine $P$ passant par $Q$ : \dotline{4cm}

\item La longueur du segment d'extrémités $T$ et $U$ : \dotline{4cm}

\item Les points $X$, $Y$ et $Z$ sont alignés : \dotline{4cm}
\end{itemize}
}

% ============================================
% EXERCICE 2 : Appartenance et positions relatives (4 points)
% ============================================
\begin{exercisebox}{\faIcon{map-marker-alt} Exercice 2 : Appartenance et positions relatives \points{4}}
\textbf{Observe la figure ci-dessous et complète :}

\exspace

\begin{center}
\begin{tikzpicture}[scale=0.8]
% Droite horizontale (d1)
\draw[thick] (-3,0) -- (5,0);
\node[right] at (5,0) {$(d_1)$};

% Droite oblique (d2)
\draw[thick] (-1,3) -- (3,-1);
\node[right] at (3,-1) {$(d_2)$};

% Points sur (d1)
\node[circle,fill=black,inner sep=1.5pt,label=below:$A$] (A) at (1,0) {};
\node[circle,fill=black,inner sep=1.5pt,label=below:$B$] (B) at (3,0) {};

% Points sur (d2) - alignés sur la droite oblique (équation: y = -x + 2)
\node[circle,fill=black,inner sep=1.5pt,label={[xshift=3mm, yshift=-2mm]$C$}] (C) at (2,0) {};
\node[circle,fill=black,inner sep=1.5pt,label={[xshift=-3mm, yshift=3mm]$D$}] (D) at (0,2) {};
\node[circle,fill=black,inner sep=1.5pt,label={[xshift=3mm, yshift=-2mm]$E$}] (E) at (1,1) {};
\end{tikzpicture}
\end{center}

\exspace

\textbf{1.} Complète avec le symbole $\in$ ou $\notin$ : \points{2}

\smallspace

$A$ \dotline{1cm} $(d_1)$ \quad ; \quad $C$ \dotline{1cm} $(d_2)$ \quad ; \quad $B$ \dotline{1cm} $(d_2)$ \quad ; \quad $D$ \dotline{1cm} $(d_1)$

\exspace

\textbf{2.} Que peut-on dire des droites $(d_1)$ et $(d_2)$ ? \points{1}

\dotline{12cm}

\exspace

\textbf{3.} Les points $C$, $D$ et $E$ sont-ils alignés ? Justifie ta réponse. \points{1}

\zonereponse{1.5cm}
\end{exercisebox}

\smallspace

% ============================================
% EXERCICE 3 : Lecture de notations (3 points)
% ============================================
\begin{exercisebox}{\faIcon{eye} Exercice 3 : Lecture de notations \points{3}}
\textbf{Donne la lecture de chaque notation :}

\smallspace

\begin{itemize}[label=$\bullet$,leftmargin=*]
\item $(AB) \parallel (CD)$ se lit : \dotline{9cm}

\vspace{0.3cm}

\item $(d) \perp (\Delta)$ se lit : \dotline{9cm}

\vspace{0.3cm}

\item Deux droites \textbf{perpendiculaires} sont deux droites qui :

\dotline{12cm}
\end{itemize}
\end{exercisebox}

\smallspace

% ============================================
% EXERCICE 4 : Constructions géométriques (4 points)
% ============================================
\begin{constructionbox}{\faIcon{drafting-compass} Exercice 4 : Constructions géométriques \points{4}}
\textbf{À l'aide de ta règle et de ton équerre, réalise les constructions suivantes sur la figure ci-dessous.\\
Code soigneusement ta figure.}

\figureespace{0.5cm}

\begin{center}
\begin{tikzpicture}[scale=1.2]
% Droite (d)
\draw[thick] (0,0) -- (8,0);
\node[right] at (8,0) {$(d)$};

% Points
\node[circle,fill=black,inner sep=1.5pt,label=above:$M$] (M) at (2,0) {};
\node[circle,fill=black,inner sep=1.5pt,label=above:$N$] (N) at (5,2) {};
\node[circle,fill=black,inner sep=1.5pt,label=above:$P$] (P) at (6.5,1.5) {};
\end{tikzpicture}
\end{center}

\exspace

\textbf{1.} Trace la perpendiculaire à $(d)$ passant par le point $M$. \points{1}

\exspace

\textbf{2.} Trace la perpendiculaire à $(d)$ passant par le point $N$. \points{1}

\exspace

\textbf{3.} Trace la parallèle à $(d)$ passant par le point $P$. \points{1}

\exspace

\textbf{4.} Code ta figure pour indiquer les angles droits et les droites parallèles. \points{1}
\end{constructionbox}

\clearpage

% ============================================
% EXERCICE 5 : Division euclidienne (4 points)
% ============================================
\begin{exercisebox}{\faIcon{divide} Exercice 5 : Division euclidienne \points{4}}

\begin{minipage}[t]{0.48\textwidth}
\textbf{1.} Pose et effectue la division euclidienne de 423 par 15. \points{2}

\vspace{0.3cm}

\zonereponse{5cm}

\end{minipage}
\hfill
\begin{minipage}[t]{0.48\textwidth}
\textbf{2.} Dans cette division euclidienne, indique : \points{1}

\vspace{0.2cm}

\begin{itemize}[label=$\bullet$,leftmargin=*,noitemsep]
\item Le dividende : \dotline{3cm}

\item Le diviseur : \dotline{3cm}

\item Le quotient : \dotline{3cm}

\item Le reste : \dotline{3cm}
\end{itemize}
\end{minipage}

\smallspace

\textbf{3.} Vérifie ton résultat en utilisant l'égalité de la division euclidienne. \points{1}

\zonereponse{2cm}

\end{exercisebox}

% Auto-évaluation
\AutoEval

\end{document}